% chapterintro
% ==============================
% 第12章 导言：12.0-chinese.md
本章的主要内容有：

\begin{itemize}
\item 使用Doxygen构建文档

\item 使用Sphinx构建文档

\item 结合Doxygen和Sphinx

\end{itemize}

文档在所有的软件项目都是有必要的：对于用户来说，了解如何获得并构建代码，并且如何有效地使用源代码或库；对于开发人员来说，文档可用来描述你源码细节，并帮助其他程序员参与其中为该项目作出贡献。本章将展示如何使用CMake构建代码文档，这里使用了两个流行的文档框架：Doxygen和Sphinx。


% --- recipe 1: 12.1-chinese.md ---
\section{12.1 使用Doxygen构建文档}

\begin{myInfoBox}
\textbf{NOTE}：此示例代码可以在 \texthref{https://github.com/dev-cafe/cmake-cookbook/tree/v1.0/chapter-12/recipe-01} 中找到，其中包含一个C++示例。该示例在CMake 3.5版(或更高版本)中是有效的，并且已经在GNU/Linux、macOS和Windows上进行过测试。
\end{myInfoBox}

Doxygen(http://www.doxygen.nl )是非常流行的源代码文档工具。可以在代码中添加文档标记作为注释，而后运行Doxygen提取这些注释，并以Doxyfile配置文件中定义的格式创建文档。Doxygen可以输出HTML、XML，甚至LaTeX或PDF。本示例将展示，如何使用CMake来构建Doxygen文档。

\mySmallsection{准备工作}

使用前几章中介绍的消息库的简化版本。目录结构如下:

\begin{codebox}[shell]
.
├── cmake
│    └── UseDoxygenDoc.cmake
├── CMakeLists.txt
├── docs
│    ├── Doxyfile.in
│    └── front_page.md
└── src
    ├── CMakeLists.txt
    ├── hello-world.cpp
    ├── Message.cpp
    └── Message.hpp
\end{codebox}

我们仍然在\texttt{src}子目录下放置源代码，并且在CMake子目录中有自定义的CMake模块。由于重点是文档，所以消除了对UUID的依赖，并简化了源代码。最大的区别是头文件中的大量代码注释：

\begin{codebox}[cpp]
#pragma once

#include <iosfwd>
#include <string>

/ * ! \file Message.hpp * /

/*! \class Message
* \brief Forwards string to screen
* \author Roberto Di Remigio
* \date 2018
* /

class Message {
public:
  /*! \brief Constructor from a string
  * \param[in] m a message
  */
  Message(const std::string &m) : message_(m) {}
  /*! \brief Constructor from a character array
  * \param[in] m a message
  */
  Message(const char * m): message_(std:: string(m)){}
  
  friend std::ostream &operator<<(std::ostream &os, Message &obj) {
    return obj.printObject(os);
  }
private:
  /*! The message to be forwarded to screen */
  std::string message_;
  /*! \brief Function to forward message to screen
  * \param[in, out] os output stream
  */
  std::ostream &printObject(std::ostream &os);
};
\end{codebox}

这些注释的格式是\texttt{/*!*/}，并包含一些Doxygen可以理解的特殊标记(参见\texthref{http://www.stack.nl/\~dimitri/Doxygen/manual/docblocks.html} )。

\mySmallsection{具体实施}

首先，来看看根目录下的\texttt{CMakeLists.txt}：

\begin{enumerate}
\item 我们声明了一个C++11项目：

\begin{codebox}[cmake]
cmake_minimum_required(VERSION 3.5 FATAL_ERROR)
project(recipe-01 LANGUAGES CXX)
set(CMAKE_CXX_STANDARD 11)
set(CMAKE_CXX_EXTENSIONS OFF)
set(CMAKE_CXX_STANDARD_REQUIRED ON)
\end{codebox}

\item 为动态库和静态库，以及可执行文件定义了输出目录：

\begin{codebox}[cmake]
include(GNUInstallDirs)
set(CMAKE_ARCHIVE_OUTPUT_DIRECTORY
    ${CMAKE_BINARY_DIR}/${CMAKE_INSTALL_LIBDIR})
set(CMAKE_LIBRARY_OUTPUT_DIRECTORY
    ${CMAKE_BINARY_DIR}/${CMAKE_INSTALL_LIBDIR})
set(CMAKE_RUNTIME_OUTPUT_DIRECTORY
    ${CMAKE_BINARY_DIR}/${CMAKE_INSTALL_BINDIR})
\end{codebox}

\item 将\texttt{cmake}子目录追加到\texttt{CMAKE\_MODULE\_PATH}。这是需要CMake找到我们的自定义模块：

\begin{codebox}[cmake]
list(APPEND CMAKE_MODULE_PATH "${CMAKE_SOURCE_DIR}/cmake")
\end{codebox}

\item \texttt{UseDoxygenDoc.cmake}自定义模块。将在后面讨论它的内容:

\begin{codebox}[cmake]
include(UseDoxygenDoc)
\end{codebox}

\item 然后添加\texttt{src}子目录:

\begin{codebox}[cmake]
add_subdirectory(src)
\end{codebox}

\end{enumerate}

\texttt{src}子目录中的\texttt{CMakeLists.txt}文件包含以下构建块:

\begin{enumerate}
\item 添加了消息库:

\begin{codebox}[cmake]
add_library(message STATIC
  Message.hpp
  Message.cpp
  )
\end{codebox}

\item 然后，声明\texttt{add\_doxygen\_doc}函数。这个函数可以理解这些参数：BUILD\_DIR、DOXY\_FILE、TARGET\_NAME和COMMENT。使用\texttt{cmake\_parse\_arguments}标准CMake命令解析这些参数：

\begin{codebox}[cmake]
function(add_doxygen_doc)
  set(options)
  set(oneValueArgs BUILD_DIR DOXY_FILE TARGET_NAME COMMENT)
  set(multiValueArgs)
  cmake_parse_arguments(DOXY_DOC
    "${options}"
    "${oneValueArgs}"
    "${multiValueArgs}"
    ${ARGN}
  )
  # ...
endfunction()
\end{codebox}

\item Doxyfile包含用于构建文档的所有Doxygen设置。一个模板\texttt{Doxyfile.in}文件作为函数参数\texttt{DOXY\_FILE}传递，并解析为\texttt{DOXY\_DOC\_DOXY\_FILE}变量。使用如下方式，配置模板文件\texttt{Doxyfile.in}:

\begin{codebox}[cmake]
configure_file(
  ${DOXY_DOC_DOXY_FILE}
  ${DOXY_DOC_BUILD_DIR}/Doxyfile
  @ONLY
  )
\end{codebox}

\item 然后，定义了一个名为\texttt{DOXY\_DOC\_TARGET\_NAME}的自定义目标，它将使用Doxyfile中的设置执行Doxygen，并在\texttt{DOXY\_DOC\_BUILD\_DIR}中输出结果:

\begin{codebox}[cmake]
add_custom_target(${DOXY_DOC_TARGET_NAME}
  COMMAND
    ${DOXYGEN_EXECUTABLE} Doxyfile
  WORKING_DIRECTORY
    ${DOXY_DOC_BUILD_DIR}
  COMMENT
    "Building ${DOXY_DOC_COMMENT} with Doxygen"
  VERBATIM
  )
\end{codebox}

\item 最后，为用户打印一条状态信息:

\begin{codebox}[cmake]
message(STATUS "Added ${DOXY_DOC_TARGET_NAME} [Doxygen] target to build documentation")
\end{codebox}

\end{enumerate}

可以像往常一样配置项目：

\begin{codebox}[shell]
$ mkdir -p build
$ cd build
$ cmake ..
$ cmake --build .
\end{codebox}

可以通过调用自定义文档目标来构建文档:

\begin{codebox}[shell]
$ cmake --build . --target docs
\end{codebox}

您将注意到构建树中出现了一个\texttt{\_build}子目录。它包含Doxygen从源文件生成的HTML文档。用浏览器打开\texttt{index.html}将显示Doxygen欢迎页面。

如果导航到类列表，例如：可以浏览Message类的文档:

\myGraphic{0.9}{images/chapter12/12-1.png}{}

\mySmallsection{工作原理}

默认情况下，CMake不支持文档构建。但是，我们可以使用\texttt{add\_custom\_target}执行任意操作。需要注意的是，需要确保构建文档所需的工具(本例中是Doxygen和Perl)在系统上可用。

此外，请注意\texttt{ UseDoxygenDoc.cmake}自定义模块只做以下工作:

\begin{itemize}
\item 执行对Doxygen和Perl可执行程序的搜索

\item 定义函数

\end{itemize}

使用\texttt{add\_doxygen\_doc}函数对文档目标进行创建。这个显式模式要优于隐式模式，我们也认为这是很好的实践方式：不要使用模块来执行类似宏(或函数)的操作。

为了限制变量定义的范围和可能出现的副作用，我们使用函数而不是宏实现了\texttt{add\_doxygen\_doc}。这种情况下，函数和宏都可以工作(并且会产生相同的结果)，但是建议优先使用函数而不是宏，除非需要修改父范围中的变量。

\begin{myInfoBox}
\textbf{NOTE}：在cmake 3.9中添加了\texttt{FindDoxygen.cmake}模块。实现了\texttt{doxygen\_add\_docs}函数，其行为与我们在本示例中给出的宏类似。要了解更多细节，请访问\texthref{https://cmake.org/cmake/help/v3.9/module/FindDoxygen.html} 查看在线文档。
\end{myInfoBox}


% --- recipe 2: 12.2-chinese.md ---
\section{12.2 使用Sphinx构建文档}

\begin{myInfoBox}
\textbf{NOTE}：此示例代码可以在 \texthref{https://github.com/dev-cafe/cmake-cookbook/tree/v1.0/chapter-12/recipe-02} 中找到，其中包含一个C++示例。该示例在CMake 3.5版(或更高版本)中是有效的，并且已经在GNU/Linux、macOS和Windows上进行过测试。
\end{myInfoBox}

Sphinx是一个Python程序，也是一个非常流行的文档系统(http://www.sphinx-doc.org )。当与Python项目一起使用时，可以为\texttt{docstring}解析源文件，并自动为函数和类生成文档页面。然而，Sphinx不仅限于Python，还可以解析reStructuredText、Markdown，并生成HTML、ePUB或PDF文档。还有在线阅读服务(https://readthedocs.org )，它提供了一种快速编写和部署文档的方法。本示例将向您展示，如何使用CMake构建Sphinx文档。

\mySmallsection{准备工作}

我们希望建立一个简单的网站，记录我们的消息库输出的信息。源码树现在看起来如下:

\begin{codebox}[shell]
.
├── cmake
│    ├── FindSphinx.cmake
│    └── UseSphinxDoc.cmake
├── CMakeLists.txt
├── docs
│    ├── conf.py.in
│    └── index.rst
└── src
    ├── CMakeLists.txt
    ├── hello-world.cpp
    ├── Message.cpp
    └── Message.hpp
\end{codebox}

\texttt{cmake}子目录中有一些自定义模块，\texttt{docs}子目录以纯文本reStructuredText格式的网站主页，\texttt{index.rst}和一个带有Sphinx的设置Python模板文件\texttt{conf.py.in}，这个模板文件可以使用\texttt{sphinx-quickstart}程序自动生成。

\mySmallsection{具体实施}

与之前的示例相比，我们将修改主\texttt{CMakeLists.txt}文件，并实现一个函数(\texttt{add\_sphinx\_doc}):

\begin{enumerate}
\item 将\texttt{cmake}文件夹附加到\texttt{CMAKE\_MODULE\_PATH}之后，我们将包括\texttt{UseSphinxDoc.cmake}自定义模块:

\begin{codebox}[cmake]
list(APPEND CMAKE_MODULE_PATH "${CMAKE_SOURCE_DIR}/cmake")
include(UseSphinxDoc)
\end{codebox}

\item \texttt{UseSphinxDoc.cmake}模块定义了\texttt{add\_sphinx\_doc}函数。我们使用关键字参数调用这个函数，以便对Sphinx文档的构建进行设置。自定义文档目标将称为\texttt{docs}:

\begin{codebox}[cmake]
add_sphinx_doc(
  SOURCE_DIR
    ${CMAKE_CURRENT_SOURCE_DIR}/docs
  BUILD_DIR
    ${CMAKE_CURRENT_BINARY_DIR}/_build
  CACHE_DIR
    ${CMAKE_CURRENT_BINARY_DIR}/_doctrees
  HTML_DIR
    ${CMAKE_CURRENT_BINARY_DIR}/sphinx_html
  CONF_FILE
    ${CMAKE_CURRENT_SOURCE_DIR}/docs/conf.py.in
  TARGET_NAME
    docs
  COMMENT
    "HTML documentation"
  )
\end{codebox}

\end{enumerate}

\texttt{UseSphinxDoc.cmake}模块遵循相同的显式方式，这样的使用方式要优于在前一个示例中的隐式方式：

\begin{enumerate}
\item 需要找到Python解释器和Sphinx可执行文件，如下:

\begin{codebox}[cmake]
find_package(PythonInterp REQUIRED)
find_package(Sphinx REQUIRED)
\end{codebox}

\item 然后，用一个值关键字参数定义\texttt{add\_sphinx\_doc}函数，并用\texttt{cmake\_parse\_arguments}解析：

\begin{codebox}[cmake]
function(add_sphinx_doc)
  set(options)
  set(oneValueArgs
    SOURCE_DIR
    BUILD_DIR
    CACHE_DIR
    HTML_DIR
    CONF_FILE
    TARGET_NAME
    COMMENT
    )
  set(multiValueArgs)
  cmake_parse_arguments(SPHINX_DOC
    "${options}"
    "${oneValueArgs}"
    "${multiValueArgs}"
    ${ARGN}
    )
  # ...
endfunction()
\end{codebox}

\item 模板文件\texttt{conf.py.in}作为\texttt{CONF\_FILE}关键字参数传递，在\texttt{SPHINX\_DOC\_BUILD\_DIR}中配置为\texttt{conf.py}：

\begin{codebox}[cmake]
configure_file(
  ${SPHINX_DOC_CONF_FILE}
  ${SPHINX_DOC_BUILD_DIR}/conf.py
  @ONLY
  )
\end{codebox}

\item 添加了一个名为\texttt{SPHINX\_DOC\_TARGET\_NAME}的自定义目标，用Sphinx来编排文档构建:

\begin{codebox}[cmake]
add_custom_target(${SPHINX_DOC_TARGET_NAME}
  COMMAND
    ${SPHINX_EXECUTABLE}
      -q
      -b html
      -c ${SPHINX_DOC_BUILD_DIR}
      -d ${SPHINX_DOC_CACHE_DIR}
      ${SPHINX_DOC_SOURCE_DIR}
      ${SPHINX_DOC_HTML_DIR}
  COMMENT
    "Building ${SPHINX_DOC_COMMENT} with Sphinx"
  VERBATIM
  )
\end{codebox}

\item 最后，打印一条状态信息:

\begin{codebox}[cmake]
message(STATUS "Added ${SPHINX_DOC_TARGET_NAME} [Sphinx] target to build documentation")
\end{codebox}

\item 配置项目并构建了文档目标:

\begin{codebox}[shell]
$ mkdir -p build
$ cd build
$ cmake ..
$ cmake --build . --target docs
\end{codebox}

\end{enumerate}

这将生成\texttt{SPHINX\_DOC\_HTML\_DIR}中的HTML文档生成树的子目录。同样，可以使用浏览器打开\texttt{index.html}，并查看文档:

\myGraphic{0.9}{images/chapter12/12-2.png}{}

\mySmallsection{工作原理}

我们利用\texttt{add\_custom\_target}的功能，可以向构建系统添加任意的构建目标。本例中，文档将使用Sphinx构建。由于Sphinx是一个可以与其他Python模块一起扩展的Python程序，所以\texttt{docs}目标将依赖于Python解释器。我们使用\texttt{find\_package}确保依赖关系。需要注意的是，\texttt{FindSphinx.cmake}模块还不是一个标准的CMake模块；它的副本包含在项目源代码中，位于\texttt{cmake}子目录下。


% --- recipe 3: 12.3-chinese.md ---
\section{12.3 结合Doxygen和Sphinx}

\begin{myInfoBox}
\textbf{NOTE}：此示例代码可以在 \texthref{https://github.com/dev-cafe/cmake-cookbook/tree/v1.0/chapter-12/recipe-03} 中找到，其中包含一个C++示例。该示例在CMake 3.5版(或更高版本)中是有效的，并且已经在GNU/Linux、macOS和Windows上进行过测试。
\end{myInfoBox}

我们有一个C++项目，因此Doxygen是生成源代码文档的理想选择。然而，我们也希望发布面向用户的文档，例如：介绍设计选择。所以我们想使用Sphinx，因为生成的HTML也可以在移动设备上查看，而且可以部署文档进行在线阅读(https://readthedocs.org )。本教程将演示如何使用Breathe插件(https://breathe.readthedocs.io )组合Doxygen和Sphinx。

\mySmallsection{准备工作}

这个示例的目录结构，类似于之前的两个示例:

\begin{codebox}[shell]
.
├── cmake
│    ├── FindPythonModule.cmake
│    ├── FindSphinx.cmake
│    └── UseBreathe.cmake
├── CMakeLists.txt
├── docs
│    ├── code-reference
│    │    ├── classes-and-functions.rst
│    │    └── message.rst
│    ├── conf.py.in
│    ├── Doxyfile.in
│    └── index.rst
└── src
    ├── CMakeLists.txt
    ├── hello-world.cpp
    ├── Message.cpp
    └── Message.hpp
\end{codebox}

\texttt{docs}子目录现在同时包含一个\texttt{ Doxyfile.in}和一个\texttt{conf.py.in}模板文件。模板文件中，分别设置了Doxygen和Sphinx。此外，还有一个\texttt{code-referenc}子目录。

\texttt{code-referenc}子目录中的文件包含Breathe指令，用来在Sphinx中包含doxygen生成的文档：

\begin{codebox}[text]
Messaging classes
=================
Message
-------
..  doxygenclass:: Message
    :project: recipe-03
    :members:
    :protected-members:
    :private-members:
\end{codebox}

这将输出Message类的文档。

\mySmallsection{具体实施}

\texttt{src}目录中的\texttt{CMakeLists.txt}文件没有改变。主\texttt{CMakeLists.txt}文件中有修改：

\begin{enumerate}
\item 包含\texttt{ UseBreathe.cmake}自定义模块：

\begin{codebox}[cmake]
list(APPEND CMAKE_MODULE_PATH "${CMAKE_SOURCE_DIR}/cmake")
include(UseBreathe)
\end{codebox}

\item 调用\texttt{add\_breathe\_doc}函数，这个函数是在自定义模块中定义的，它接受关键字参数，来设置Doxygen和Sphinx：

\begin{codebox}[cmake]
add_breathe_doc(
  SOURCE_DIR
    ${CMAKE_CURRENT_SOURCE_DIR}/docs
  BUILD_DIR
    ${CMAKE_CURRENT_BINARY_DIR}/_build
  CACHE_DIR
    ${CMAKE_CURRENT_BINARY_DIR}/_doctrees
  HTML_DIR
    ${CMAKE_CURRENT_BINARY_DIR}/html
  DOXY_FILE
    ${CMAKE_CURRENT_SOURCE_DIR}/docs/Doxyfile.in
  CONF_FILE
    ${CMAKE_CURRENT_SOURCE_DIR}/docs/conf.py.in
  TARGET_NAME
    docs
  COMMENT
    "HTML documentation"
  )
\end{codebox}

\end{enumerate}

让我们看一下\texttt{UseBreatheDoc.cmake}模块，其遵循了与我们在前两个示例中描述的显式模式。具体描述如下:

\begin{enumerate}
\item 文档生成依赖于Doxygen:

\begin{codebox}[cmake]
find_package(Doxygen REQUIRED)
find_package(Perl REQUIRED)
\end{codebox}

\item 还依赖于Python解释器和Sphinx:

\begin{codebox}[cmake]
find_package(PythonInterp REQUIRED)
find_package(Sphinx REQUIRED)
\end{codebox}

\item 此外，还必须找到breathe的Python模块。这里，我们使用\texttt{FindPythonModule.cmake}模块:

\begin{codebox}[cmake]
include(FindPythonModule)
find_python_module(breathe REQUIRED)
\end{codebox}

\item 定义了\texttt{add\_breathe\_doc}函数，这个函数有一个单值关键字参数，我们将使用\texttt{cmake\_parse\_arguments}命令解析它:

\begin{codebox}[cmake]
function(add_breathe_doc)
  set(options)
  set(oneValueArgs
    SOURCE_DIR
    BUILD_DIR
    CACHE_DIR
    HTML_DIR
    DOXY_FILE
    CONF_FILE
    TARGET_NAME
    COMMENT
    )
  set(multiValueArgs)
  cmake_parse_arguments(BREATHE_DOC
    "${options}"
    "${oneValueArgs}"
    "${multiValueArgs}"
    ${ARGN}
    )
  # ...
endfunction()
\end{codebox}

\item \texttt{BREATHE\_DOC\_CONF\_FILE}中的Sphinx模板文件，会通过\texttt{conf.py}配置到的\texttt{BREATHE\_DOC\_BUILD\_DIR}目录下：

\begin{codebox}[cmake]
configure_file(
  ${BREATHE_DOC_CONF_FILE}
  ${BREATHE_DOC_BUILD_DIR}/conf.py
  @ONLY
  )
\end{codebox}

\item 相应地，Doxygen的\texttt{BREATHE\_DOC\_DOXY\_FILE}模板文件配置为\texttt{BREATHE\_DOC\_BUILD\_DIR}中的Doxyfile:

\begin{codebox}[cmake]
configure_file(
  ${BREATHE_DOC_DOXY_FILE}
  ${BREATHE_DOC_BUILD_DIR}/Doxyfile
  @ONLY
  )
\end{codebox}

\item 添加\texttt{BREATHE\_DOC\_TARGET\_NAME}自定义目标。注意，只有Sphinx在运行时，对Doxygen的调用才发生在\texttt{BREATHE\_DOC\_SPHINX\_FILE}中:

\begin{codebox}[cmake]
add_custom_target(${BREATHE_DOC_TARGET_NAME}
  COMMAND
    ${SPHINX_EXECUTABLE}
      -q
      -b html
      -c ${BREATHE_DOC_BUILD_DIR}
      -d ${BREATHE_DOC_CACHE_DIR}
      ${BREATHE_DOC_SOURCE_DIR}
      ${BREATHE_DOC_HTML_DIR}
  COMMENT
    "Building ${BREATHE_DOC_TARGET_NAME} documentation with Breathe, Sphinx and Doxygen"
  VERBATIM
  )
\end{codebox}

\item 最后，打印一条状态信息:

\begin{codebox}[cmake]
message(STATUS "Added ${BREATHE_DOC_TARGET_NAME} [Breathe+Sphinx+Doxygen] target to build documentation")
\end{codebox}

\item 配置完成后，构建文档:

\begin{codebox}[shell]
$ mkdir -p build
$ cd build
$ cmake ..
$ cmake --build . --target docs
\end{codebox}

\end{enumerate}

该文档将在\texttt{BREATHE\_DOC\_HTML\_DIR}子目录中可用。启动浏览器打开\texttt{index.html}文件后，可以导航到Message类的文档:

\myGraphic{0.9}{images/chapter12/12-3.png}{}

\mySmallsection{工作原理}

尽管在声明定制的\texttt{BREATHE\_DOC\_TARGET\_NAME}目标时只调用了Sphinx，但这里Doxygen和Sphinx都在运行。这要感谢Sphinx的\texttt{conf.py}文件中的以下设置:

\begin{codebox}[python]
def run_doxygen(folder):
  """Run the doxygen make command in the designated folder"""

  try:
      retcode = subprocess.call("cd {}; doxygen".format(folder), shell=True)
    if retcode < 0:
      sys.stderr.write(
      "doxygen terminated by signal {}".format(-retcode))
  except OSError as e:
    sys.stderr.write("doxygen execution failed: {}".format(e))


def setup(app):
  run_doxygen('@BREATHE_DOC_BUILD_DIR@')
\end{codebox}

Doxygen将生成XML输出，Breathe插件将能够与所选择的Sphinx文档样式一致的形式，呈现XML输出。


